\chapter{Rangkaian Listrik dan Daya}\label{ch:g11:circuits-and-power}

Sentakkan sebuah sakelar: sebuah lampu menyala, karena sebuah pembangkit
yang berjarak berkilo-kilometer mendorong muatan mengelilingi sebuah
untai tembaga yang melewati dindingmu. Tahun-tahun sebelumnya sudah
memetakan rangkaiannya --- arus, \omterm{def:g11:circuits-and-power:voltage}{tegangan}, \omterm{prop:g11:circuits-and-power:ohm}{hukum Ohm}. Bab ini
membukanya kembali dengan buku besar
\cref{ch:g10:energy-conservation}: apa yang ditagih meterannya, mengapa
kawat yang \omterm{def:g11:lenses-and-eye:lens}{tipis} berpijar, dan mengapa jaringannya bekerja pada
\qty{400}{kV}.

\section{Muatan yang bergerak}

\begin{definition}[Arus dan coulomb]\label{def:g11:circuits-and-power:current}
\emph{Arus listrik}\index{arus listrik} adalah aliran muatan yang
teratur (elektron, di dalam logam). Kuat arusnya $I$, dalam
\emph{ampere}\index{ampere} (\unit{A}), adalah muatan yang melewati
penampang kawatnya tiap sekon: arus tetap $I$ yang mengalir selama
waktu $t$ mengangkut muatan
\[ q = I\,t , \]
dalam \emph{coulomb}\index{coulomb} (\unit{C}), dengan $\qty{1}{C} =
\qty{1}{A.s}$. Muatan dasarnya adalah $e = \qty{1.6e-19}{C}$. Arus
diukur dengan amperemeter, yang dipasang seri.
\end{definition}

\begin{example}[Petir lawan telepon]\label{ex:g11:circuits-and-power:coulomb}
Sebuah sambaran petir membawa sekitar $\qty{3.0e4}{A}$ selama kira-kira
$\qty{1.0e-4}{s}$:
$q = 3.0\times10^{4} \times 10^{-4} = \qty{3.0}{C}$. Teleponmu, pada
\qty{2.0}{A}, memindahkan muatan yang sama tiap $1.5$ sekon.
\end{example}

\begin{definition}[Tegangan dan volt]\label{def:g11:circuits-and-power:voltage}
\emph{Tegangan}\index{tegangan} $U$ antara dua titik sebuah rangkaian,
dalam \emph{volt}\index{volt} (\unit{V}), adalah \omterm{def:g10:energy-conservation:energy}{energi} yang
dipertukarkan tiap \omterm{def:g10:orders-of-magnitude:unit}{satuan} muatan yang berjalan di antara keduanya:
$\qty{1}{V} = \qty{1}{J/C}$, sehingga muatan $q$ yang menyeberangi
sebuah alat pada tegangan $U$ mempertukarkan \omterm{def:g10:energy-conservation:energy}{energi} $E = qU$ dengannya.
Voltmeter, yang dipasang paralel, mengukurnya.
\end{definition}

\begin{proposition}[Hukum rangkaian]\label{prop:g11:circuits-and-power:laws}
Pada rangkaian yang bekerja mantap:
\begin{enumerate}
  \item \emph{hukum titik cabang:} arus yang tiba pada sebuah titik
        cabang berjumlah sama dengan arus yang meninggalkannya;
  \item \emph{hukum seri:} \omterm{def:g11:circuits-and-power:voltage}{tegangan} sepanjang sebuah \omterm{def:g10:relative-motion:trajectory}{lintasan}
        dijumlahkan: $U_{AC} = U_{AB} + U_{BC}$;
  \item \emph{hukum paralel:} dua cabang yang bertemu pada dua titik
        yang sama membawa \omterm{def:g11:circuits-and-power:voltage}{tegangan} yang sama.
\end{enumerate}
\end{proposition}

\begin{proof}
Pada keadaan mantap muatan tidak menumpuk di mana pun, sehingga apa yang
mengalir masuk ke sebuah titik cabang mengalir keluar juga. \omterm{def:g10:energy-conservation:energy}{Energi} tiap
\omterm{def:g11:circuits-and-power:current}{coulomb} bersifat menjumlah sepanjang \omterm{def:g10:relative-motion:trajectory}{lintasan} dan hanya bergantung pada
kedua ujungnya --- bukan pada cabang yang ditempuh di antaranya.
\end{proof}

\section{Hukum Ohm dan jaringan hambatan}

\begin{definition}[Hambatan]\label{def:g11:circuits-and-power:resistance}
\emph{Hambatan}\index{hambatan} sebuah penghantar adalah nisbah
$R = U/I$ antara \omterm{def:g11:circuits-and-power:voltage}{tegangan} padanya dan arus yang melaluinya, yang diukur
dalam \emph{ohm}\index{ohm} (\unit{\ohm}): $\qty{1}{\ohm} = \qty{1}{V/A}$.
\end{definition}

\begin{proposition}[Hukum Ohm]\label{prop:g11:circuits-and-power:ohm}
\index{hukum Ohm}
Untuk penghantar logam yang dijaga pada suhu tetap, $R$ bernilai tetap:
\omterm{def:g11:circuits-and-power:voltage}{tegangan} dan arusnya sebanding,
\[ U = R\,I . \]
\end{proposition}
\admitted

\begin{remark}\label{rem:g11:circuits-and-power:ohm-limits}
Ini hukum percobaan, bukan hukum semesta: dioda, kawat pijar lampu
pemanas atau kulitmu sendiri tidak menaatinya. Mengapa logam menaatinya
diturunkan pada jilid perguruan tinggi.
\end{remark}

\begin{proposition}[Hambatan seri dan paralel]\label{prop:g11:circuits-and-power:networks}
\index{hambatan pengganti}
Dua \omterm{def:g11:circuits-and-power:resistance}{hambatan} yang dipasang seri setara dengan satu \omterm{def:g11:circuits-and-power:resistance}{hambatan}
$R_{\text{s}} =
R_1 + R_2$; dua \omterm{def:g11:circuits-and-power:resistance}{hambatan} paralel setara dengan satu \omterm{def:g11:circuits-and-power:resistance}{hambatan}
$R_{\text{p}}$ yang diberikan oleh
\[
\frac{1}{R_{\text{p}}} = \frac{1}{R_1} + \frac{1}{R_2} ,
\qquad\text{yaitu}\quad
R_{\text{p}} = \frac{R_1 R_2}{R_1 + R_2} < \min(R_1, R_2) .
\]
\end{proposition}

\begin{proof}
Seri: arus $I$ yang sama menyeberangi keduanya dan \omterm{def:g11:circuits-and-power:voltage}{tegangannya}
dijumlahkan (\cref{prop:g11:circuits-and-power:laws}), sehingga
$U = (R_1 + R_2)\,I$. Paralel: keduanya berada pada $U$ yang sama dan
arusnya dijumlahkan, sehingga $I = \frac{U}{R_1} + \frac{U}{R_2}$.
\end{proof}

\begin{omfigure}
\begin{circuitikz}[font=\small]
  \draw (0,0) to[battery1, l=$U$] (0,2.2)
        to[ammeter] (1.9,2.2)
        to[R=$R_1$] (3.7,2.2)
        to[R=$R_2$] (5.5,2.2) -- (5.5,0) -- (0,0);
  \draw (3.7,2.2) -- (3.7,3.3) to[voltmeter] (5.5,3.3) -- (5.5,2.2);
\end{circuitikz}\qquad
\begin{circuitikz}[font=\small]
  \draw (0,0) to[battery1, l=$U$] (0,2.2)
        to[ammeter, l=$I$] (2.2,2.2);
  \draw (2.2,2.2) to[R=$R_1$, i=$I_1$] (2.2,0) -- (0,0);
  \draw (2.2,2.2) -- (4.0,2.2) to[R=$R_2$, i=$I_2$] (4.0,0) -- (2.2,0);
\end{circuitikz}

{\small Seri (kiri): satu arus, \omterm{def:g11:circuits-and-power:voltage}{tegangannya} dijumlahkan. Paralel
(kanan): satu \omterm{def:g11:circuits-and-power:voltage}{tegangan}, arusnya dijumlahkan.}
\end{omfigure}

\begin{example}[Hambatan pengganti]\label{ex:g11:circuits-and-power:networks}
$R_1 = \qty{100}{\ohm}$ dan $R_2 = \qty{150}{\ohm}$: dipasang seri,
\qty{250}{\ohm}; dipasang paralel, $\frac{100 \times 150}{250} =
\qty{60}{\ohm}$ --- lebih kecil daripada keduanya: arusnya memperoleh
jalan kedua.
\end{example}

\section{Energi dan daya di dalam rangkaian}

\begin{proposition}[Daya listrik]\label{prop:g11:circuits-and-power:power}
\index{daya listrik}
Sebuah alat pada tegangan $U$ yang dilalui arus $I$ menerima \omterm{def:g10:energy-conservation:energy}{energi}
dengan laju
\[ P = U\,I , \qquad\text{yaitu}\quad E = U\,I\,t \ \text{dalam waktu } t . \]
\end{proposition}

\begin{proof}
Dalam waktu $t$ muatan $q = It$ menyeberangi alat itu
(\cref{def:g11:circuits-and-power:current}) dan tiap \omterm{def:g11:circuits-and-power:current}{coulomb}
menyerahkan $U$ \omterm{def:g10:energy-conservation:energy}{joule} (\cref{def:g11:circuits-and-power:voltage}):
$E = qU = UIt$.
\end{proof}

\begin{proposition}[Daya pada sebuah hambatan]\label{prop:g11:circuits-and-power:ri2}
Sebuah \omterm{def:g11:circuits-and-power:resistance}{hambatan} $R$ yang dialiri arus $I$ membuang \omterm{def:g10:energy-conservation:power}{daya}
\[ P = R\,I^2 = \frac{U^2}{R} . \]
\end{proposition}

\begin{proof}
Sisipkan \omterm{prop:g11:circuits-and-power:ohm}{hukum Ohm} ke dalam $P = UI$:
$P = (RI)I = RI^2 = U(U/R) = U^2/R$.
\end{proof}

\begin{example}[Membaca papan nama]\label{ex:g11:circuits-and-power:nameplate}
Sebuah ketel bertanda ``\qty{230}{V} -- \qty{2200}{W}'' menarik
$I = P/U =
2200/230 = \qty{9.6}{A}$ dan berhambatan
$R = U^2/P = 230^2/2200 =
\qty{24}{\ohm}$. Menyala \qty{150}{s} ia
memakai $E = 2200 \times 150 =
\qty{3.3e5}{J} \approx \qty{0.09}{kW.h}$
--- sekitar dua sen.
\end{example}

\begin{remark}[Joule dan kilowatt-jam]\label{rem:g11:circuits-and-power:kwh}
Meterannya menghitung \omterm{def:g10:energy-conservation:kwh}{kilowatt-jam}
(\cref{def:g10:energy-conservation:kwh}):
$\qty{1}{kW.h} = \qty{3.6e6}{J}$. Joule adalah \omterm{def:g10:orders-of-magnitude:unit}{satuan} milik fisikawan;
\omterm{def:g10:energy-conservation:kwh}{kilowatt-jam} yang menghargai tagihannya.
\end{remark}

\begin{omfigure}
\begin{tikzpicture}[font=\small, x=1.5cm, y=0.42cm]
  \draw[-Stealth] (0,0) -- (4.6,0) node[right] {$P$};
  \foreach \x/\t in {1/\qty{10}{W}, 2/\qty{100}{W}, 3/\qty{1}{kW}, 4/\qty{10}{kW}}
    {\draw (\x,-0.3) -- (\x,0.3); \node[below, font=\footnotesize] at (\x,-0.35) {\t};}
  \fill[omDef!70] (0,1) rectangle (1.0,1.75);
  \node[right] at (1.05,1.4) {lampu LED (\qty{10}{W})};
  \fill[omDef!70] (0,2.35) rectangle (1.78,3.1);
  \node[right] at (1.83,2.75) {komputer jinjing (\qty{60}{W})};
  \fill[omThm!70] (0,3.7) rectangle (2.95,4.45);
  \node[right] at (3.0,4.1) {pemanggang roti (\qty{900}{W})};
  \fill[omThm!70] (0,5.05) rectangle (3.34,5.8);
  \node[right] at (3.39,5.45) {ketel (\qty{2.2}{kW})};
  \fill[omThm!70] (0,6.4) rectangle (3.95,7.15);
  \node[right] at (4.0,6.8) {pemanas air mandi (\qty{9}{kW})};
\end{tikzpicture}

{\small \omterm{def:g10:energy-conservation:power}{Daya} peralatan pada sumbu logaritma: segala sesuatu yang
memanaskan (merah) hidup sepuluh kali lipat di atas segala sesuatu yang
menghitung atau bercahaya (biru).}
\end{omfigure}

\section{Efek Joule: dari pemanggang roti sampai menara listrik}

\begin{definition}[Efek Joule]\label{def:g11:circuits-and-power:joule}
Pembuangan \omterm{def:g10:energy-conservation:power}{daya} $P = RI^2$ yang mengubah \omterm{def:g10:energy-conservation:forms}{energi listrik} menjadi kalor
pada setiap \omterm{def:g11:circuits-and-power:resistance}{hambatan} disebut \emph{efek Joule}\index{efek Joule}.
Pemanas ruangan, ketel, pemanggang roti dan bohlam model lama memang
dirancang sebagai \omterm{def:g11:circuits-and-power:resistance}{hambatan}; sebuah \emph{sekring}\index{sekring} adalah
kawat \omterm{def:g11:lenses-and-eye:lens}{tipis} yang ditera agar melebur --- lalu memutus rangkaiannya ---
ketika arusnya melampaui nilai teranya.
\end{definition}

\begin{example}[Mengapa menara listrik bertegangan tinggi]\label{ex:g11:circuits-and-power:transmission}
Sebuah desa memerlukan $P = \qty{100}{kW}$ melalui saluran berhambatan
$R_\ell = \qty{5.0}{\ohm}$. Diserahkan pada $U = \qty{500}{V}$:
$I = P/U =
\qty{200}{A}$ dan salurannya memboroskan
$R_\ell I^2 = \qty{200}{kW}$ --- dua kali yang diserahkannya: mustahil.
Pada $U = \qty{20}{kV}$: $I = \qty{5.0}{A}$ dan
$R_\ell I^2 = \qty{125}{W}$. Arusnya empat puluh kali lebih kecil,
rugi \omterm{def:g10:energy-conservation:power}{dayanya} $1600$ kali lebih kecil.
\end{example}

\begin{omfigure}
\begin{circuitikz}[font=\small]
  \draw (0,0) to[sV, l=pembangkit] (0,2)
        to[R, l_=$R_\ell$, i=$I$] (3.4,2) -- (4.8,2)
        to[lamp, l=desa] (4.8,0) -- (0,0);
  \node[font=\footnotesize, omThm] at (1.7,2.75) {rugi \omterm{def:g10:energy-conservation:power}{daya} $R_\ell I^2$};
\end{circuitikz}

{\small Penyerahan \omterm{def:g10:energy-conservation:power}{daya} pada
\cref{ex:g11:circuits-and-power:transmission}: menaikkan $U$ pada $P = UI$ yang tetap menyusutkan $I$, dan rugi \omterm{def:g10:energy-conservation:power}{dayanya} turun bersama $I^2$.}
\end{omfigure}

\begin{remark}[Aturan besi jaringan listrik]\label{rem:g11:circuits-and-power:grid}
Pada \omterm{def:g10:energy-conservation:power}{daya} terserah $P = UI$ yang tetap, rugi saluran
$R_\ell I^2 = R_\ell
P^2/U^2$ turun sebagai \emph{kuadrat} \omterm{def:g11:circuits-and-power:voltage}{tegangan}
pengangkutannya. Karena itu: naikkan \omterm{def:g11:circuits-and-power:voltage}{tegangannya} (\qty{400}{kV}) untuk
menempuh perjalanan, lalu turunkan (\qty{230}{V}) untuk dipakai
sehari-hari.
\end{remark}

\begin{remark}[Keselamatan rumah tangga]\label{rem:g11:circuits-and-power:safety}
Saluran \qty{230}{V} yang disekring pada \qty{16}{A} memasok paling
banyak $230 \times 16
\approx \qty{3.7}{kW}$: dua ketel sekaligus
meleburkan \omterm{def:g11:circuits-and-power:joule}{sekringnya} --- memang dirancang sebagai mata rantai yang
paling \omterm{def:g11:lenses-and-eye:lens}{tipis} dan paling lemah, jauh lebih murah daripada kebakaran di
dalam dinding.
\end{remark}

\section{Baterai yang sebenarnya: gaya gerak listrik}

\begin{definition}[GGL dan hambatan dalam]\label{def:g11:circuits-and-power:emf}
Sebuah baterai mengubah \omterm{def:g10:energy-conservation:forms}{energi kimia} menjadi \omterm{def:g10:energy-conservation:forms}{energi listrik}. \emph{Gaya
gerak listrik}\index{gaya gerak listrik} (GGL) $\mathcal{E}$-nya, dalam
\omterm{def:g11:circuits-and-power:voltage}{volt}, adalah \omterm{def:g10:energy-conservation:energy}{energi} yang diserahkannya kepada tiap \omterm{def:g11:circuits-and-power:current}{coulomb} yang
menyeberanginya; \emph{hambatan dalam}\index{hambatan dalam} $r$-nya
menampung \omterm{def:g11:circuits-and-power:joule}{efek Joule} di dalam kimianya sendiri. Modelnya: sumber ideal
$\mathcal{E}$ yang dipasang seri dengan $r$ ($\mathcal{E}$ yang
berlekuk: $E$ adalah \omterm{def:g10:energy-conservation:energy}{energi}).
\end{definition}

\begin{proposition}[Tegangan jepit]\label{prop:g11:circuits-and-power:terminal}
Sebuah baterai yang menyerahkan arus $I$ menunjukkan pada kutubnya
\[ U = \mathcal{E} - r\,I . \]
\end{proposition}

\begin{proof}
Kimianya memasok $\mathcal{E}I$, \omterm{def:g11:circuits-and-power:emf}{hambatan dalamnya} memakan $rI^2$, dan
rangkaiannya menerima $UI$: \omterm{thm:g10:energy-conservation:conservation}{kekekalan energi}
(\cref{ch:g10:energy-conservation}) memberikan
$UI = \mathcal{E}I - rI^2$.
\end{proof}

\begin{example}[Baterai yang dibebani]\label{ex:g11:circuits-and-power:discharge}
Sebuah baterai ($\mathcal{E} = \qty{9.0}{V}$, $r = \qty{1.0}{\ohm}$)
memasok bohlam berhambatan $R = \qty{8.0}{\ohm}$. Untainya menaati
$\mathcal{E} =
rI + RI$, sehingga $I = \mathcal{E}/(R + r) = \qty{1.0}{A}$
dan $U = 9.0 -
1.0 \times 1.0 = \qty{8.0}{V}$: bohlamnya menerima
\qty{8.0}{W} sementara baterainya memboroskan \qty{1.0}{W} untuk
memanaskan dirinya sendiri.
\end{example}

\begin{omfigure}
\begin{tikzpicture}
\begin{axis}[omaxis, width=8cm, height=5cm,
    xlabel={$I$ (\unit{A})}, ylabel={$U$ (\unit{V})},
    xmin=0, xmax=4.6, ymin=0, ymax=10.8,
    xtick={0,1,2,3,4}, ytick={0,3,6,9}]
  \addplot[omThm, thick, domain=0:4.3] {9 - x};
  \addplot[omDef, only marks, mark=*, mark size=1.5pt]
    coordinates {(0.5,8.5) (1.0,8.0) (2.0,7.0) (3.0,6.0) (4.0,5.0)};
  \node[font=\small, anchor=west] at (axis cs:0.2,10.1)
    {titik potong $\mathcal{E} = \qty{9.0}{V}$};
\end{axis}
\end{tikzpicture}

{\small Ciri pengosongan baterai pada
\cref{ex:g11:circuits-and-power:discharge}: titik yang terukur berada
pada garis $U = \mathcal{E} - rI$, dengan kemiringan
$-r = \qty{-1.0}{V/A}$.}
\end{omfigure}

\begin{remark}[Penerima]\label{rem:g11:circuits-and-power:receivers}
Sebuah motor atau baterai yang sedang diisi menjalankan pengubahannya
secara terbalik: ia menyerap pada $U = \mathcal{E}' + r'I$, mengubah
$\mathcal{E}'I$ menjadi gerak atau kimia dan kehilangan $r'I^2$. Sumber
mendorong, penerima mendorong balik.
\end{remark}

\begin{example}[Efisiensi sebuah baterai]\label{ex:g11:circuits-and-power:efficiency}
\omterm{def:g10:energy-conservation:efficiency}{Efisiensi} (\cref{def:g10:energy-conservation:efficiency}) pengosongan di
atas: yang bergunanya $UI = \qty{8.0}{W}$ dari $\mathcal{E}I =
\qty{9.0}{W}$, sehingga $\eta = U/\mathcal{E} = 8/9 \approx 0.89$ ---
dan nilainya turun ketika arusnya naik: baterai yang bekerja keras
menjadi panas.
\end{example}

\section{Latihan}

\begin{exercise}[$\star$]\label{exo:g11:circuits-and-power:1}
Sebuah baterai telepon diisi dengan arus tetap \qty{2.0}{A} selama
\qty{90}{min}. Berapa muatan yang melewati kabelnya, dalam \omterm{def:g11:circuits-and-power:current}{coulomb}?
Setara berapa muatan dasar ($e = \qty{1.6e-19}{C}$)?
\end{exercise}

\begin{exercise}[$\star$]\label{exo:g11:circuits-and-power:2}
Sebuah \omterm{def:g11:circuits-and-power:resistance}{hambatan} \qty{220}{\ohm} dipasang pada \qty{12}{V}: berapa arus
yang mengalir? Sebuah elemen pemanas menarik \qty{0.50}{A} pada
\qty{230}{V}: berapa \omterm{def:g11:circuits-and-power:resistance}{hambatannya}?
\end{exercise}

\begin{exercise}[$\star$]\label{exo:g11:circuits-and-power:3}
Hitunglah \omterm{prop:g11:circuits-and-power:networks}{hambatan pengganti} $R_1 = \qty{120}{\ohm}$ dan
$R_2 =
\qty{60}{\ohm}$ ketika dipasang seri, lalu ketika dipasang
paralel. Yang mana yang lebih kecil daripada keduanya, dan mengapa?
\end{exercise}

\begin{exercise}[$\star$]\label{exo:g11:circuits-and-power:4}
Sebuah setrika perjalanan menarik \qty{8.7}{A} pada \qty{230}{V}.
Hitunglah \omterm{def:g10:energy-conservation:power}{dayanya}, lalu \omterm{def:g10:energy-conservation:energy}{energi} yang dipakainya dalam \qty{4.0}{min},
dalam \omterm{def:g10:energy-conservation:energy}{joule} dan dalam \omterm{def:g10:energy-conservation:kwh}{kilowatt-jam}.
\end{exercise}

\begin{exercise}[$\star$]\label{exo:g11:circuits-and-power:5}
Sebuah pengisi \omterm{def:g10:energy-conservation:power}{daya} komputer jinjing menyerahkan \qty{60}{W} selama
\qty{5.0}{h} tiap hari: berapa \omterm{def:g10:energy-conservation:energy}{energi} hariannya dalam \omterm{def:g10:energy-conservation:kwh}{kilowatt-jam} dan
\omterm{def:g10:energy-conservation:energy}{joule}, dan berapa biayanya pada $0.25$ euro tiap \unit{kW.h}?
\end{exercise}

\begin{exercise}[$\star\star$]\label{exo:g11:circuits-and-power:6}
Sebuah pemanas listrik ditera \qty{1200}{W} pada \qty{230}{V}.
\begin{enumerate}
  \item Hitunglah \omterm{def:g11:circuits-and-power:resistance}{hambatannya} dan arus yang ditariknya.
  \item Jika dicolokkan (dengan \omterm{def:g11:circuits-and-power:resistance}{hambatan} yang sama) ke pasokan
        \qty{115}{V}, berapa \omterm{def:g10:energy-conservation:power}{daya} yang diserahkannya? Mengapa bukan
        separuhnya?
\end{enumerate}
\end{exercise}

\begin{exercise}[$\star\star$]\label{exo:g11:circuits-and-power:7}
$R_1 = \qty{40}{\ohm}$ dan $R_2 = \qty{20}{\ohm}$ dipasang seri pada
\qty{12}{V}. Hitunglah arusnya, \omterm{def:g11:circuits-and-power:voltage}{tegangan} pada tiap \omterm{def:g11:circuits-and-power:resistance}{hambatannya}, dan
\omterm{def:g10:energy-conservation:power}{daya} yang dibuang masing-masing; periksalah bahwa \omterm{def:g10:energy-conservation:power}{dayanya} berjumlah
$UI$.
\end{exercise}

\begin{exercise}[$\star\star$]\label{exo:g11:circuits-and-power:8}
$R_1 = \qty{60}{\ohm}$ dan $R_2 = \qty{30}{\ohm}$ dipasang paralel pada
\qty{12}{V}. Hitunglah arus tiap cabangnya, arus seluruhnya, dan
\omterm{prop:g11:circuits-and-power:networks}{hambatan penggantinya} dengan dua cara.
\end{exercise}

\begin{exercise}[$\star\star$]\label{exo:g11:circuits-and-power:9}
Sebuah baterai menyerahkan $U = \qty{5.5}{V}$ pada $I = \qty{0.50}{A}$,
dan $U = \qty{4.0}{V}$ pada $I = \qty{2.0}{A}$. Tentukan GGL-nya
$\mathcal{E}$ dan \omterm{def:g11:circuits-and-power:emf}{hambatan dalamnya} $r$, lalu arus hubung singkatnya
$\mathcal{E}/r$.
\end{exercise}

\begin{exercise}[$\star\star$]\label{exo:g11:circuits-and-power:10}
Sebuah saluran dapur \qty{230}{V} dilindungi \omterm{def:g11:circuits-and-power:joule}{sekring} \qty{16}{A}. Sebuah
ketel \qty{2.2}{kW} dan pemanggang roti \qty{1.0}{kW} menyala bersamaan:
apakah \omterm{def:g11:circuits-and-power:joule}{sekringnya} bertahan? Seseorang menambahkan pemanas \qty{1.5}{kW}:
tunjukkan bahwa \omterm{def:g11:circuits-and-power:joule}{sekringnya} putus.
\end{exercise}

\begin{exercise}[$\star\star$]\label{exo:g11:circuits-and-power:11}
Sebuah bengkel memerlukan \qty{10}{kW} melalui kabel penghantar
berhambatan \qty{2.0}{\ohm}. Hitunglah arus dan rugi Joule-nya jika
\omterm{def:g10:energy-conservation:power}{dayanya} diserahkan pada \qty{200}{V}, lalu pada \qty{4.0}{kV}.
Bandingkan kedua rugi \omterm{def:g10:energy-conservation:power}{dayanya} dan jelaskan nisbahnya.
\end{exercise}

\begin{exercise}[$\star\star\star$]\label{exo:g11:circuits-and-power:12}
Dengan memakai ketiganya sekaligus, yaitu tiga \omterm{def:g11:circuits-and-power:resistance}{hambatan} \qty{60}{\ohm}
yang serupa, daftarkan setiap \omterm{prop:g11:circuits-and-power:networks}{hambatan pengganti} yang berbeda (seri,
paralel, campuran); hitunglah masing-masing.
\end{exercise}

\begin{exercise}[$\star\star\star$]\label{exo:g11:circuits-and-power:13}
Sebuah baterai ($\mathcal{E} = \qty{4.5}{V}$, $r = \qty{1.5}{\ohm}$)
memasok \omterm{def:g11:circuits-and-power:resistance}{hambatan} $R = \qty{6.0}{\ohm}$.
\begin{enumerate}
  \item Hitunglah $I$, $U$, \omterm{def:g10:energy-conservation:power}{daya} berguna pada $R$, \omterm{def:g10:energy-conservation:power}{daya} yang hilang
        pada $r$, dan \omterm{def:g10:energy-conservation:efficiency}{efisiensinya}.
  \item Tunjukkan bahwa \omterm{def:g10:energy-conservation:power}{daya} yang diterima $R$ yang dapat diubah-ubah,
        $P = \mathcal{E}^2 R/(R+r)^2$, bernilai terbesar ketika
        $R = r$ (petunjuk: $(R+r)^2/R = (R-r)^2/R + 4r$), lalu
        hitunglah nilai terbesarnya.
\end{enumerate}
\end{exercise}

\begin{exercise}[$\star\star\star$]\label{exo:g11:circuits-and-power:14}
Sebuah ketel \qty{2.2}{kW} berefisiensi $0.90$ membawa \qty{1.0}{kg} air
dari \qty{20}{\degreeCelsius} sampai mendidih (memanaskan air memakan
\qty{4180}{J} tiap \omterm{def:g10:orders-of-magnitude:unit}{kilogram} tiap derajat). Hitunglah kalor yang
diperlukan, \omterm{def:g10:energy-conservation:forms}{energi listrik} yang ditarik, lama pemanasannya, dan biayanya
pada $0.25$ euro tiap \unit{kW.h}.
\end{exercise}

\begin{exercise}[$\star\star\star$]\label{exo:g11:circuits-and-power:15}
Sebuah aki mobil: $\mathcal{E} = \qty{12.7}{V}$,
$r = \qty{0.020}{\ohm}$, kapasitasnya \qty{60}{A.h}.
\begin{enumerate}
  \item Ubahlah kapasitasnya ke \omterm{def:g11:circuits-and-power:current}{coulomb}, lalu taksirlah \omterm{def:g10:energy-conservation:energy}{energi} yang
        tersimpan ($E \approx \mathcal{E} q$) dalam \omterm{def:g10:energy-conservation:energy}{joule} dan
        \omterm{def:g10:energy-conservation:kwh}{kilowatt-jam}.
  \item Motor starternya menarik \qty{150}{A}: hitunglah \omterm{def:g11:circuits-and-power:voltage}{tegangan}
        jepitnya, \omterm{def:g10:energy-conservation:power}{daya} yang diserahkannya, dan \omterm{def:g10:energy-conservation:power}{daya} yang hilang di
        dalamnya.
  \item Jelaskan mengapa lampu depannya meredup selagi mesinnya
        distarter.
\end{enumerate}
\end{exercise}

\section{Soal: Dari bendungan sampai pemanggang roti}

\begin{problem}[{Soal akhir pekan --- perjalanan seribu kilometer satu
kilowatt-jam, dari air yang jatuh sampai roti yang mencokelat, beserta
seluruh tagihan perjalanannya}]
\label{pb:g11:circuits-and-power:1}
Sebuah bendungan gunung menahan air \qty{200}{m} di atas turbinnya, yang
menelan \qty{50}{m^3} --- \qty{5.0e4}{kg} --- tiap sekon; turbin dan
alternatornya mengubah dengan \omterm{def:g10:energy-conservation:efficiency}{efisiensi} $0.90$. Melalui saluran
berhambatan seluruhnya $R_\ell = \qty{32}{\ohm}$ (\qty{1000}{km} pergi
dan kembali) pembangkit itu memasok sebuah kota, tempat pemanggang roti
\qty{900}{W} milikmu menunggu pada \qty{230}{V}. Alternatornya
mengeluarkan \qty{20}{kV}; transformator (\omterm{def:g10:energy-conservation:efficiency}{efisiensi} $0.99$ masing-masing)
mengubah \omterm{def:g11:circuits-and-power:voltage}{tegangannya} sesuka hati. Listriknya dijual pada $0.25$ euro
tiap \unit{kW.h}; $g = \qty{9.81}{N/kg}$.

\medskip
\noindent\textbf{Bagian I --- Di bendungan.}
\begin{enumerate}
  \item Berapa \omterm{def:g10:energy-conservation:energy}{energi} potensial yang dilepaskan air yang jatuh itu tiap
        sekon? Simpulkan \omterm{def:g10:energy-conservation:power}{daya} yang tersedia bagi turbinnya.
  \item Berapa \omterm{prop:g11:circuits-and-power:power}{daya listrik} yang meninggalkan alternatornya?
  \item Pada \qty{20}{kV} langsung dari alternatornya, berapa arus yang
        mengalir?
  \item Hitunglah rugi Joule $R_\ell I^2$ pada arus itu lalu bandingkan
        dengan keluaran pembangkitnya. Simpulkan.
  \item \omterm{def:g10:energy-conservation:power}{Daya} kotanya $P = UI$ ditetapkan oleh permintaan. Satu-satunya
        cara apa untuk menyusutkan $I$, dan alat yang mana yang
        melakukannya?
\end{enumerate}

\medskip
\noindent\textbf{Bagian II --- Di saluran.}
\begin{enumerate}[resume]
  \item Sebuah transformator, yang untuk pertanyaan ini dianggap tanpa
        rugi \omterm{def:g10:energy-conservation:power}{daya}, menaikkan \qty{88}{MW} itu ke \qty{400}{kV}: berapa
        arus yang mengalir sekarang?
  \item Hitunglah rugi salurannya yang baru, dalam megawatt.
  \item Berapa bagian keluaran pembangkitnya yang hilang di salurannya?
  \item Tunjukkan bahwa bagian yang hilang itu sama dengan
        $R_\ell P/U^2$, lalu periksalah bahwa hal itu menjelaskan faktor
        antara pertanyaan 4 dan 7.
  \item Hitunglah jatuh tegangan $R_\ell I$ sepanjang salurannya pada
        \qty{400}{kV}. Apakah nilainya sesuai dengan pertanyaan 8?
\end{enumerate}

\medskip
\noindent\textbf{Bagian III --- Di dapur.}
\begin{enumerate}[resume]
  \item Ubahlah \qty{1}{kW.h} ke \omterm{def:g10:energy-conservation:energy}{joule}; apa yang dikalikan meterannya?
  \item Hitunglah arus dan \omterm{def:g11:circuits-and-power:resistance}{hambatan} pemanggang rotinya pada \qty{230}{V}.
  \item Tiga menit memanggang: berapa \omterm{def:g10:energy-conservation:energy}{energinya} dalam \omterm{def:g10:energy-conservation:energy}{joule} dan
        \omterm{def:g10:energy-conservation:kwh}{kilowatt-jam}, dan berapa harganya?
  \item Pemanggang roti itu membuang $RI^2$: jelaskan dalam satu kalimat
        mengapa rumus yang sama menjadi rugi \omterm{def:g10:energy-conservation:power}{daya} di salurannya dan
        menjadi barang jadinya di dapur.
  \item Pemanggang roti itu berbagi \omterm{def:g11:circuits-and-power:joule}{sekring} \qty{16}{A} dengan ketel
        \qty{2.2}{kW}: apakah \omterm{def:g11:circuits-and-power:joule}{sekringnya} bertahan? Sebuah pemanas
        \qty{2.0}{kW} bergabung: bagaimana sekarang?
\end{enumerate}

\medskip
\noindent\textbf{Bagian IV --- Auditnya.}
\begin{enumerate}[resume]
  \item Kabel rumahnya sendiri, yang berhambatan \qty{0.05}{\ohm},
        membawa \qty{13.5}{A} pada pertanyaan 15: hitunglah rugi \omterm{def:g10:energy-conservation:power}{dayanya}
        dan bagiannya dari kira-kira \qty{3.1}{kW} yang diserahkan.
  \item Rangkaikan \omterm{def:g10:energy-conservation:efficiency}{efisiensinya} (turbin--alternator, penaik \omterm{def:g11:circuits-and-power:voltage}{tegangan},
        saluran, penurun \omterm{def:g11:circuits-and-power:voltage}{tegangan}, kabel rumah) menjadi satu \omterm{def:g10:energy-conservation:efficiency}{efisiensi}
        jaringan seluruhnya.
  \item Berapa \omterm{def:g10:orders-of-magnitude:unit}{kilogram} (dan berapa \omterm{def:g10:orders-of-magnitude:unit}{meter} kubik) air bendungan yang
        harus jatuh agar \emph{satu} \omterm{def:g10:energy-conservation:kwh}{kilowatt-jam} sampai ke pemanggang
        rotimu?
  \item Hitunglah bagian yang hilang $R_\ell P/U^2$ pada pertanyaan 9
        untuk $U = \qty{20}{kV}$. Apa artinya nilai yang lebih besar
        daripada $1$?
  \item Puncak ceritanya --- \emph{harga sebuah jarak}: dalam dua
        kalimat, berapa harga satu \omterm{def:g10:energy-conservation:kwh}{kilowatt-jam} yang memanggang roti itu
        bagi gunungnya, berapa bagiannya yang mati di perjalanan, dan
        satu muslihat yang mana yang membuatnya mungkin?
\end{enumerate}
\end{problem}
